\section{问题三：光伏预报展开与滚动修正模型}

\subsection{预报的时间对齐}

附件 3 每行是一次预报，“预报 $h$ 小时”指发布时刻起第 $h$ 个小时，即区间 $(\text{发布时刻}+h-1,\ \text{发布时刻}+h]$。例如 6:00 发布的“预报 1 小时”对应当天 6:00--7:00，也就是第 36 个十分钟区间。因此发布时刻 $r$ 的预报展开后正好覆盖区间 $[6r,144)$：
\begin{equation}
PV^{\mathrm{fc}(r)}_{6r+i}=\bigl[\text{附件3 在 }r\text{ 时发布的第 }(i/6+1)\text{ 小时值}\bigr]\,\Delta t,
\qquad i=0,\dots,144-6r-1 .
\label{eq:p3expand}
\end{equation}
即在该小时的 6 个十分钟区间内取常值。

滚动修正在电力市场与微网调度中已有广泛应用\cite{morales2014,wang2020}。

\subsection{滚动修正模型}

\textbf{0:00 初始计划。}用附件 3 的 0:00 预报作为光伏预测、式~\eqref{eq:fc} 的 7 日均值作为负载预测，加上安全裕度式~\eqref{eq:margin}，得到初始计划 $(x^{\text{plan}}_t,c_t,d_t)$。

\textbf{$r\in\{6,12,18\}$ 的调整。}在 $r$ 时刻，区间 $[0,6r)$ 已经执行完毕、不可更改；只对剩余区间 $[6r,144)$ 重新求解线性规划，起点为当前储电量，收尾仍回到当天起始电量。由于引入了更准的预报，剩余时段的光伏预测被替换为 $PV^{\mathrm{fc}(r)}$。

\subsection{偏差结算}

题面给出的两条规则可以合并。设 $\Delta_t=x_t^{\text{adj}}-x_t^{\text{plan}}$。

\textbf{增购}（$\Delta_t>0$）：计划部分按 $p_t$ 计费，超出部分按 $1.5p_t$ 购入，
\begin{equation}
p_tx^{\text{plan}}_t+1.5p_t\Delta_t=p_tx^{\text{adj}}_t+0.5p_t\Delta_t .
\end{equation}

\textbf{减购}（$\Delta_t<0$）：调整后的电量按 $p_t$ 计费，被削减的部分按 $0.5p_t$ 承担违约费，
\begin{equation}
p_tx^{\text{adj}}_t+0.5p_t\bigl(x^{\text{plan}}_t-x^{\text{adj}}_t\bigr)
=p_tx^{\text{adj}}_t+0.5p_t|\Delta_t| .
\end{equation}

两式在各自的分支上完全一致，因此可以统一写成对称的偏差结算式：
\begin{equation}
C_{\text{购电}}=\sum_{t}p_tx^{\text{adj}}_t
+\frac12\sum_t p_t\bigl|x^{\text{adj}}_t-x^{\text{plan}}_t\bigr| .
\label{eq:deviation}
\end{equation}
总费用为 $C=C_{\text{购电}}+5\sum_t p_t E_t$。式~\eqref{eq:deviation} 中的第二项是凸的分段线性函数，用一对非负变量 $u_t,v_t\ge0$ 且 $x^{\text{adj}}_t-u_t+v_t=x^{\text{plan}}_t$ 线性化后仍然是一个标准线性规划。

\subsection{预报精度分析}

表 7 把三种光伏预测放在\textbf{同一时段}上比较平均绝对误差，避免因覆盖窗口不同而产生偏差。

\threelinetable{表 7\quad 三种光伏预测在同一时段上的平均绝对误差（kWh/区间）}{lrrr}
{时段 & 前 7 日均值 & 只用 0:00 预报 & 滚动使用最新预报}
{0:00--6:00   & 0.7  & 6.6   & 6.6 \\
 6:00--12:00  & 52.5 & 114.2 & 98.4 \\
 12:00--18:00 & 48.7 & 123.4 & 92.2 \\
 18:00--24:00 & 0.4  & 3.2   & 3.2 \\
 全天整体     & 25.6 & 61.9  & ---}

\begin{figure}[H]
  \centering
  \includegraphics[width=0.86\textwidth]{forecast_accuracy.pdf}
  \caption{同一时段下三种光伏预测的平均绝对误差对比}
  \label{fig:fc}
\end{figure}

两个结论值得注意。第一，\textbf{滚动使用最新预报确实降低了误差}：在 6:00--12:00 段由 $114.2$ 降到 $98.4\ \mathrm{kWh}$，在 12:00--18:00 段由 $123.4$ 降到 $92.2\ \mathrm{kWh}$，降幅分别为 $13.8\%$ 和 $25.3\%$。第二，\textbf{附件 3 的预报本身并不比“前 7 天的实际均值”更准}——全天整体误差 $61.9$ 对 $25.6$，几乎是后者的 2.4 倍。这解释了为什么问题三的总费用高于问题二（同样的安全裕度机制下，预报误差更大）。这也说明，如果允许同时使用两路信息，把附件 3 的预报与历史均值做加权融合还有进一步降本的空间。

\subsection{滚动修正的结果与结论}

\threelinetable{表 8\quad 问题三全年结果对比（2025-02-01 至 2025-12-31）}{lrr}
{指标 & 引入 6/12/18 点预报 & 只用 0:00 预报}
{0:00 计划购电量 / kWh & 25012959.00 & 25012959.00 \\
 调整后购电量 / kWh & 24611319.48 & 25012959.00 \\
 购电费 / 元 & 16130122.56 & 16416828.97 \\
 偏差违约费 / 元 & 454289.21 & 0.00 \\
 紧急购电量 / kWh & 268200.65 & 366664.04 \\
 紧急购电费 / 元 & 1021169.31 & 1389336.28 \\
 \textbf{总购电费用 / 元} & \textbf{17605581.08} & \textbf{17806165.24}}

\textbf{结论：值得引入后续预报。}滚动方案比只用 0:00 预报的方案净省
$17806165.24-17605581.08=\mathbf{200584.16}$ 元/年。分解来看，紧急购电费下降 $36.8$ 万元，而偏差违约费只有 $45.4$ 万元、调整后购电费本身又下降了 $28.7$ 万元，三者相抵后仍净赚 $20.06$ 万元。

另一个值得注意的现象是\textbf{调整方向以减购为主}：调整后购电量 $24611319.48\,\mathrm{kWh}$ 小于 0:00 计划的 $25012959.00\,\mathrm{kWh}$。这说明 0:00 计划在安全裕度的作用下偏保守，后续预报主要起“纠正过度保守”的作用——在确认当天不会缺电之后，把多余的保险退掉（代价是 $0.5$ 倍的违约费）。

\threelinetable{表 9\quad 问题三指定日期的结果}{lrrr}
{日期 & 0:00 计划 / kWh & 调整后 / kWh & 紧急购电 / kWh}
{2025-03-20 & 68068.16 & 66797.62 & 1976.85 \\
 2025-06-21 & 73080.10 & 73578.04 & 0.00 \\
 2025-09-23 & 70044.31 & 66709.12 & 1160.65 \\
 2025-12-21 & 96561.06 & 94489.77 & 1461.53}

四天对应的总费用分别为 $51528.93$、$49997.26$、$50931.20$、$66322.47$ 元。与问题二同期相比，3 月 20 日费用上升（$45119.57\to51528.93$ 元），但 12 月 21 日的紧急购电由 $649.18$ 升到 $1461.53\,\mathrm{kWh}$ 而购电费下降，最终总费用由 $63386.35$ 元升到 $66322.47$ 元——同样是附件 3 预报精度不如历史均值所致。
